
\chapter{Proposta}\label{cap:proposta}

Este capítulo descreve o sistema a ser implementado, o ambiente do experimento,
as configurações dos mecanismos do gRPC que serão comparadas e a forma como o
desempenho, a tolerância a falhas, a segurança e o ponto de ruptura serão
avaliados.

\section{Sistema a ser implementado}

O sistema é um \textit{marketplace} de comércio eletrônico em Go, composto por
cinco serviços (inventário, busca, compra, pagamento simulado e usuários)
atrás de um \textit{gateway}. A comunicação entre os serviços usa gRPC e
\textit{Protocol Buffers}, e o Envoy atua como proxy e balanceador de
carga entre eles, com suas políticas próprias de \textit{retry}, \textit{timeout}
e \textit{circuit breaking} desabilitadas, de modo que o efeito observado seja
atribuído aos mecanismos do gRPC. Os mecanismos são configurados nos clientes
gRPC dos serviços, e as chamadas são unárias.

Três fluxos de negócio compõem a carga: gerenciar o inventário, buscar itens e
comprar. A compra é o fluxo central, pois é uma transação distribuída (bloqueio
  do carrinho, atualização de estoque e pagamento, com \textit{rollback} em caso
de falha) que atravessa vários serviços. Ela gera a cadeia de chamadas mais
profunda e é o ponto natural para a injeção de falhas.

A persistência usa três tecnologias. O PostgreSQL armazena os dados transacionais de compras,
usuários e pagamentos, e o MongoDB armazena o inventário e as solicitações de compra
(\textit{tickets}). Em vez de uma instância por serviço, cada tecnologia tem uma única instância
compartilhada, com esquemas ou coleções separados por serviço, de modo que um serviço não acessa os
dados de outro. Essa escolha reduz o consumo de memória na VPS e mantém a comunicação entre os
serviços restrita ao gRPC. O Redis é usado de duas formas: como cache de leitura do
inventário, que reduz a carga do MongoDB nas operações de busca, e como \textit{lock} distribuído
com TTL, que bloqueia os itens do carrinho durante a compra e evita que um bloqueio fique preso
caso o serviço falhe antes de liberá-lo.

\section{Infraestrutura e ambiente}

O experimento roda em duas VPS da Contabo com Ubuntu 24.04, cada uma com 4 vCPUs e 8~GB de memória.
Uma executa o sistema, e a outra executa a geração de carga (k6) e a observabilidade (Prometheus e
Grafana), de modo que o gerador de carga não compete por recursos com o sistema testado. As
métricas do sistema são coletadas pelo Prometheus, e as métricas do lado do cliente, pelo k6. As
versões de Go, gRPC-Go e Envoy serão fixadas na implementação do sistema. Para garantir a
reprodutibilidade, o código e os arquivos de configuração serão publicados em um repositório no GitHub.

Os serviços, os bancos de dados e o Envoy são orquestrados com o Kubernetes, na distribuição leve
k3s, em um único nó na VPS do sistema. Cada componente executa em um pod com limites de
CPU e memória definidos nos manifestos, e o Kubernetes também é a base da injeção de falhas
(Seção~\ref{sec:falhas}). Os manifestos fazem parte do repositório publicado.

\section{Configurações experimentais}

A hipótese central do trabalho é que os mecanismos nativos do gRPC contribuem
para a resiliência e a disponibilidade do sistema, mas com custos que variam
conforme a carga e o tipo de falha. Para avaliá-la, o sistema é executado em
cinco configurações progressivas dos mecanismos (Tabela~\ref{tab:configuracoes}).
O \textit{keepalive} é também variado nos cenários de falha de conexão. Como o
gRPC não permite que um \textit{retry} e o \textit{hedging} regam a mesma
chamada \cite{gRFC-A6}, a configuração completa aplica cada um a métodos
diferentes.

\begin{table}[htb]
  \centering
  \caption{Configurações comparadas}\label{tab:configuracoes}
  \begin{tabular}{lcccc}
    \hline
    Configuração & \textit{Deadlines} & \textit{Retries} & \textit{Hedging} & \textit{Keepalive} \\
    \hline
    (a) sem mecanismos & não & não & não & não \\
    (b) \textit{deadlines} & sim & não & não & não \\
    (c) \textit{retries}   & sim & sim & não & não \\
    (d) \textit{hedging}   & sim & não & sim & não \\
    (e) completa           & sim & sim & sim & sim \\
    \hline
  \end{tabular}
\end{table}

Os \textit{deadlines} valem o dobro do tempo médio de resposta de cada serviço
em operação normal. Nos \textit{retries}, o número máximo de tentativas é 4 (a
original mais três retentativas), com o \textit{backoff} exponencial do gRPC.

Os \textit{retries} são acionados apenas pelo código \texttt{UNAVAILABLE}, que indica falha
transitória, com \textit{backoff} inicial de 100~ms, máximo de 1~s e multiplicador 2. O
\textit{retry throttling} é usado nas configurações (c) e (e), com \texttt{maxTokens} igual a 10 e
\texttt{tokenRatio} igual a 0,1. No \textit{hedging}, o atraso entre as chamadas é aproximadamente o
p95 do tempo de resposta do serviço em operação normal, o número máximo de tentativas é 3 e o
código \texttt{UNAVAILABLE} é tratado como falha não fatal, de modo que as demais chamadas
continuam em andamento. No \textit{keepalive}, o cliente envia um \textit{ping} a cada 10~s e
encerra a conexão se não houver confirmação em 3~s, inclusive sem chamadas ativas, e o servidor é
configurado para aceitar \textit{pings} com intervalo mínimo compatível. Como os \textit{retries}
e o \textit{hedging} podem duplicar o efeito de operações de escrita, as escritas da compra e do
pagamento carregam uma chave de idempotência, e o serviço descarta a execução repetida de uma
mesma chave. Assim, esses mecanismos podem ser aplicados também a elas. Esses valores são
iniciais e podem ser ajustados na implementação, sendo registrados no repositório.

\section{Avaliação de desempenho}

A carga é gerada com o k6, que suporta gRPC e o fluxo composto da compra, com
90\% de operações de leitura e 10\% de escrita. A carga segue um modelo fechado
(usuários concorrentes), tanto constante quanto crescente. Cada execução tem 60
segundos de \textit{warm-up} e 10 minutos de medição. As métricas são a latência
(p50, p90 e p99) e a vazão, coletadas pelo k6 e pelo Prometheus, e o uso de CPU
e memória, coletado pelo Prometheus.

\section{Tolerância a falhas e disponibilidade}\label{sec:falhas}

As falhas são injetadas nos serviços periféricos, com o sistema sob carga
crescente, para verificar se o sistema se mantém disponível diante de uma falha
parcial. Os tipos de falha são escolhidos para exercitar cada mecanismo:
latência artificial (\textit{deadlines}), erros de indisponibilidade
(\textit{retries}), lentidão de cauda (\textit{hedging}) e queda de conexão
(\textit{keepalive}).

As falhas são injetadas com o Chaos Mesh, que opera sobre os pods do Kubernetes. A
latência artificial e a lentidão de cauda são produzidas com o \texttt{NetworkChaos}, que adiciona
atraso e variação (\textit{jitter}) ao tráfego do serviço alvo. A indisponibilidade e a queda de
conexão são produzidas com o \texttt{PodChaos} (falha ou reinício do pod) e com o
\texttt{NetworkChaos} (perda de pacotes e partição de rede). Cada experimento define a duração e o
momento da injeção em um manifesto versionado, o que permite repeti-lo nas cinco configurações.

O MTTR é medido da seguinte forma. A falha é considerada iniciada quando a taxa
de requisições com erro ultrapassa um limiar, definido a partir das métricas
obtidas em momentos de pico de carga, e é considerada resolvida quando o sistema
volta a responder dentro desse limiar por 10 segundos consecutivos. A
disponibilidade é a proporção de requisições bem-sucedidas em relação ao total.

\section{Análise de segurança}

A segurança é avaliada em três frentes. A primeira é o uso do TLS do gRPC na
comunicação entre os serviços. A segunda é a superfície de ataque (portas
expostas e \textit{endpoints} acessíveis) e as vulnerabilidades conhecidas nas
dependências, buscadas com o Docker Scout e o \texttt{govulncheck}. A terceira é
o efeito de amplificação de carga causado pelos \textit{retries} e pelo
\textit{hedging}, medido em cada configuração pela razão entre as chamadas
recebidas pelos serviços e as chamadas originadas pelos clientes. Este trabalho
mede esse efeito, mas não executa ataques.

Para organizar a análise, é construído um modelo de ameaças com o STRIDE, que classifica as
ameaças em falsificação de identidade, adulteração, repúdio, divulgação de informação, negação de
serviço e elevação de privilégio. O modelo é aplicado aos componentes do sistema e às fronteiras de
confiança entre o \textit{gateway}, os serviços, os bancos de dados e o Redis, e cada ameaça é
associada ao controle que a mitiga ou à avaliação deste trabalho que a cobre. A negação de serviço
recebe atenção especial, pois os \textit{retries} e o \textit{hedging} podem ampliar a carga sobre
os serviços.

\section{Ponto de ruptura}

O ponto de ruptura é buscado nas configurações (a) e (e), para comparar o
sistema sem mecanismos com o sistema com todos eles. A carga é aumentada em
rampa, de 100 requisições por segundo a cada 10 segundos, até que o sistema
deixe de atender aos critérios: p99 acima de 500~ms ou taxa de erro acima de
1\%. O gargalo é identificado pelas métricas de CPU, memória, rede e latência
do Prometheus e do k6.

\section{Análise dos dados e cronograma}

Os resultados são analisados por médias e intervalos de confiança e apresentados
em tabelas e gráficos. Entre as ameaças à validade, o sistema é uma aplicação
que apenas se parece com um sistema de produção: seus requisitos têm
semelhanças com os reais, mas há muitos outros requisitos e implementações em
ambientes reais. O cronograma é apresentado na Tabela~\ref{tab:cronograma}.

\begin{table}[htb]
  \centering
  \caption{Cronograma}\label{tab:cronograma}
  \begin{tabular}{ll}
    \hline
    Data & Etapa \\
    \hline
    02/12/2026 & Apresentação do TCC01 \\
    28/02/2027 & Implementação do sistema \\
    15/03/2027 & Instrumentação do sistema \\
    30/03/2027 & Experimentos \\
    13/04/2027 & Análise dos resultados e escrita do TCC02 \\
    19/04/2027 & Apresentação do TCC02 \\
    \hline
  \end{tabular}
\end{table}
